\documentclass{article}
\usepackage{amsmath,amssymb,amsthm,algorithm,algorithmic}
\usepackage{bm}
\usepackage{hyperref}
\newtheorem{theorem}{Theorem}
\newtheorem{lemma}{Lemma}
\newtheorem{assumption}{Assumption}
\newcommand{\E}{\mathbb{E}}
\newcommand{\R}{\mathbb{R}}

\begin{document}

\section*{AMSGrad}

\section*{Mathematical Formulation}
Let $f:\R^{d}\rightarrow\R$ be the (possibly non‑convex) objective.  
The algorithm maintains two exponential moving averages
\[
\begin{aligned}
\bm{m}_{t} &:= \beta_{1}\bm{m}_{t-1}+(1-\beta_{1})\bm{g}_{t},\\[2mm]
\bm{v}_{t} &:= \beta_{2}\bm{v}_{t-1}+(1-\beta_{2})\bm{g}_{t}\odot\bm{g}_{t},
\end{aligned}
\tag{1}
\]
where $\bm{g}_{t}=\nabla f(\bm{w}_{t})$ (or an unbiased stochastic gradient) and
$\odot$ denotes element‑wise product.  
Bias‑corrected moments are
\[
\hat{\bm{m}}_{t}:=\frac{\bm{m}_{t}}{1-\beta_{1}^{t}},\qquad
\hat{\bm{v}}_{t}:=\frac{\bm{v}_{t}}{1-\beta_{2}^{t}}.
\tag{2}
\]
AMSGrad enforces a monotone non‑increasing stepsize by defining a
\emph{max‑tracked} second‑moment
\[
\tilde{\bm{v}}_{t}:=\max\!\bigl(\tilde{\bm{v}}_{t-1},\hat{\bm{v}}_{t}\bigr),\qquad 
\tilde{\bm{v}}_{0}:=\mathbf{0}.
\tag{3}
\]
The parameter update is
\[
\boxed{\;
\bm{w}_{t+1}= \bm{w}_{t}
-\alpha\,
\frac{\hat{\bm{m}}_{t}}{\sqrt{\tilde{\bm{v}}_{t}}+\varepsilon\mathbf{1}}\;,
\tag{4}
\]
with learning rate $\alpha>0$, decay rates $\beta_{1},\beta_{2}\in[0,1)$ and
stability constant $\varepsilon>0$.

\section*{Pseudocode}
\begin{algorithm}[H]
\caption{AMSGrad}
\begin{algorithmic}[1]
\STATE \textbf{Input:} $\bm{w}_{1}\in\R^{d}$, stepsize $\alpha$, 
$\beta_{1},\beta_{2}\in[0,1)$, $\varepsilon>0$.
\STATE Initialize $\bm{m}_{0}\gets\mathbf{0},\;
\bm{v}_{0}\gets\mathbf{0},\;
\tilde{\bm{v}}_{0}\gets\mathbf{0}$.
\FOR{$t=1,2,\dots,T$}
\STATE Compute (stochastic) gradient $\bm{g}_{t}=\nabla f(\bm{w}_{t})$.
\STATE $\bm{m}_{t}\gets\beta_{1}\bm{m}_{t-1}+(1-\beta_{1})\bm{g}_{t}$.
\STATE $\bm{v}_{t}\gets\beta_{2}\bm{v}_{t-1}+(1-\beta_{2})\bm{g}_{t}\odot\bm{g}_{t}$.
\STATE $\hat{\bm{m}}_{t}\gets\bm{m}_{t}/(1-\beta_{1}^{t})$,
        $\hat{\bm{v}}_{t}\gets\bm{v}_{t}/(1-\beta_{2}^{t})$.
\STATE $\tilde{\bm{v}}_{t}\gets\max(\tilde{\bm{v}}_{t-1},\hat{\bm{v}}_{t})$
        (element‑wise maximum).
\STATE $\bm{w}_{t+1}\gets\bm{w}_{t}
          -\alpha\;\hat{\bm{m}}_{t}\oslash(\sqrt{\tilde{\bm{v}}_{t}}+\varepsilon\mathbf{1})$,
        where $\oslash$ is element‑wise division.
\ENDFOR
\STATE \textbf{Output:} $\bm{w}_{T+1}$ (or the averaged iterate).
\end{algorithmic}
\end{algorithm}

\section*{Dry Run Example}
Consider the one‑dimensional quadratic $f(w)=(w-3)^{2}$.
The true gradient is $g(w)=2(w-3)$.  
Choose $\alpha=0.1$, $\beta_{1}=0.9$, $\beta_{2}=0.999$, $\varepsilon=10^{-8}$,
and initialise $w_{1}=0$, $m_{0}=v_{0}=\tilde v_{0}=0$.

\begin{center}
\begin{tabular}{c|c|c|c|c|c}
$t$ & $w_{t}$ & $g_{t}=2(w_{t}-3)$ & $m_{t}$ & $v_{t}$ & $w_{t+1}$\\\hline
1 & $0$ & $-6$ &
$m_{1}=0.1\cdot(-6)=-0.6$ &
$v_{1}=0.001\cdot36=0.036$ &
$\hat m_{1}=-6$, $\hat v_{1}=36$, $\tilde v_{1}=36$,\\
 & & & & & $w_{2}=0-0.1\frac{-6}{\sqrt{36}}=0+0.1=0.1$\\\hline
2 & $0.1$ & $-5.8$ &
$m_{2}=0.9(-0.6)+0.1(-5.8)=-1.14$ &
$v_{2}=0.999\cdot0.036+0.001\cdot33.64\approx0.0696$ &
$\hat m_{2}= -5.8$, $\hat v_{2}=33.64$, $\tilde v_{2}=36$ (max unchanged),\\
 & & & & & $w_{3}=0.1-0.1\frac{-5.8}{\sqrt{36}}=0.1+0.0967\approx0.1967$\\\hline
3 & $0.1967$ & $-5.6066$ &
$m_{3}=0.9(-1.14)+0.1(-5.6066)=-1.726$ &
$v_{3}=0.999\cdot0.0696+0.001\cdot31.44\approx0.0995$ &
$\hat m_{3}\approx-5.6$, $\hat v_{3}\approx31.44$, $\tilde v_{3}=36$,\\
 & & & & & $w_{4}=0.1967-0.1\frac{-5.6}{\sqrt{36}}\approx0.1967+0.0933=0.29$\\
\end{tabular}
\end{center}
Objective values:
$f(0)=9$, $f(0.1)=8.41$, $f(0.1967)\approx7.86$, $f(0.29)\approx7.35$,
showing descent.

\newpage

\section*{Assumptions}
\begin{assumption}[Smoothness]\label{ass:smooth}
$f$ is $L$-smooth, i.e.
\[
\forall\bm{x},\bm{y}\in\R^{d}:\quad 
f(\bm{y})\le f(\bm{x})+\langle\nabla f(\bm{x}),\bm{y}-\bm{x}\rangle
+\frac{L}{2}\|\bm{y}-\bm{x}\|^{2}. 
\tag{5}
\]
\end{assumption}

\begin{assumption}[Bounded Gradient]\label{ass:bound}
There exists $G>0$ such that $\|\nabla f(\bm{w})\|\le G$ for all $\bm{w}$.
\tag{6}
\end{assumption}

\begin{assumption}[Lower Bounded Objective]\label{ass:lb}
$f^{\star}:=\inf_{\bm{w}}f(\bm{w})>-\infty$.
\tag{7}
\end{assumption}

\begin{assumption}[Hyper‑parameter regime]\label{ass:params}
The decay rates satisfy $0\le\beta_{1}<\sqrt{\beta_{2}}<1$ and the stepsize
$\alpha>0$ is constant. Moreover $\varepsilon>0$.
\tag{8}
\end{assumption}

\section*{Main Convergence Theorem}
\begin{theorem}[Non‑convex Convergence of AMSGrad]\label{thm:main}
Under Assumptions \ref{ass:smooth}–\ref{ass:params}, let $\{\bm{w}_{t}\}_{t=1}^{T}$
be the iterates produced by AMSGrad with deterministic gradients.
Define $C_{1}:=\frac{G^{2}}{(1-\beta_{1})^{2}}$, $C_{2}:=\frac{L}{2}$.
Then
\[
\frac{1}{T}\sum_{t=1}^{T}\E\!\left[\|\nabla f(\bm{w}_{t})\|^{2}\right]
\;\le\;
\frac{f(\bm{w}_{1})-f^{\star}}{\alpha (1-\beta_{1})\sqrt{1-\beta_{2}}\,T}
\;+\;
\frac{C_{1}C_{2}\,\alpha}{\sqrt{1-\beta_{2}}\,T}
\;+\;
\frac{C_{1}}{T}\sum_{i=1}^{d}\frac{1}{\sqrt{\tilde v_{T,i}}}.
\tag{9}
\]
Consequently, if $\alpha=O(1/\sqrt{T})$ the right–hand side is
$O\!\bigl(\frac{\log T}{\sqrt{T}}\bigr)$; in particular
\[
\min_{1\le t\le T}\E\!\left[\|\nabla f(\bm{w}_{t})\|^{2}\right]
=O\!\Bigl(\frac{\log T}{\sqrt{T}}\Bigr).
\]
\end{theorem}

\section*{Theoretical Properties}
\begin{itemize}
\item \textbf{Stability.} The monotone max operation \eqref{3} guarantees that the denominator
$\sqrt{\tilde{\bm v}_{t}}$ never decreases, preventing the stepsize from exploding.
\item \textbf{Hyper‑parameter Sensitivity.}
Condition $0\le\beta_{1}<\sqrt{\beta_{2}}$ (Assumption \ref{ass:params})
ensures that the momentum term does not dominate the adaptivity
and is essential for the bound in \eqref{9}.
\item \textbf{Comparison to Adam.}
If the max operation is omitted, the same derivation yields a term
$\frac{1}{\sqrt{\hat{\bm v}_{t}}}$ that can become arbitrarily small,
invalidating the telescoping argument; AMSGrad’s extra term
removes this pathological behaviour, restoring the $O(\log T/\sqrt{T})$ rate.
\end{itemize}

\section*{Discussion}
The theorem shows that AMSGrad attains the same worst‑case convergence
rate as the original Adam analysis for non‑convex smooth objectives,
while the monotonicity of $\tilde{\bm v}_{t}$ yields a provably
stable stepsize. The rate matches the optimal $O(1/\sqrt{T})$ up to a
logarithmic factor originating from the element‑wise max.

\newpage

\section*{Preliminaries (Definitions and Lemmas)}
\begin{lemma}[Young’s inequality]\label{lem:young}
For any $\bm{a},\bm{b}\in\R^{d}$ and $\gamma>0$,
\[
\langle\bm{a},\bm{b}\rangle\le\frac{\|\bm{a}\|^{2}}{2\gamma}
+\frac{\gamma\|\bm{b}\|^{2}}{2}.
\tag{10}
\]
\end{lemma}
\begin{proof}
Immediate from $(\sqrt{\gamma}\bm{a}-\bm{b}/\sqrt{\gamma})^{2}\ge0$.
\end{proof}

\begin{lemma}[Bound on bias‑corrected first moment]\label{lem:m-bound}
For all $t$, $\|\hat{\bm{m}}_{t}\|\le\frac{G}{1-\beta_{1}}$.
\tag{11}
\]
\begin{proof}
From the recursion $\bm{m}_{t}= \beta_{1}\bm{m}_{t-1}+(1-\beta_{1})\bm{g}_{t}$ and
$\|\bm{g}_{t}\|\le G$, an induction yields
$\|\bm{m}_{t}\|\le G(1-\beta_{1}^{t})/(1-\beta_{1})$.
Dividing by $1-\beta_{1}^{t}$ gives \eqref{11}.
\end{proof}

\begin{lemma}[Monotonicity of $\tilde{\bm v}_{t}$]\label{lem:vt-monotone}
$\tilde{\bm v}_{t}\ge\tilde{\bm v}_{t-1}$ element‑wise for all $t$.
\tag{12}
\]
\begin{proof}
Definition \eqref{3} is a component‑wise maximum with the previous
value, hence the inequality holds by construction.
\end{proof}

\section*{Main Proof (Step‑by‑Step Derivation)}
We prove Theorem \ref{thm:main}.  
For readability, denote $\bm{w}_{t+1}=\bm{w}_{t}-\alpha\bm{d}_{t}$ where
\[
\bm{d}_{t}:=
\frac{\hat{\bm{m}}_{t}}{\sqrt{\tilde{\bm v}_{t}}+\varepsilon\mathbf{1}}.
\tag{13}
\]

\noindent\textbf{[Step 1]} Apply $L$‑smoothness \eqref{5} with
$\bm{y}=\bm{w}_{t+1}$, $\bm{x}=\bm{w}_{t}$:
\[
f(\bm{w}_{t+1})\le f(\bm{w}_{t})
-\alpha\langle\nabla f(\bm{w}_{t}),\bm{d}_{t}\rangle
+\frac{L\alpha^{2}}{2}\|\bm{d}_{t}\|^{2}.
\tag{14}
\]

\noindent\textbf{[Step 2]} Expand the inner product using the definition
of $\bm{d}_{t}$ and apply Young’s inequality \eqref{lem:young} with
$\gamma=\frac{1}{\sqrt{1-\beta_{2}}}$:
\[
\begin{aligned}
\langle\nabla f(\bm{w}_{t}),\bm{d}_{t}\rangle
&= \sum_{i=1}^{d}
\frac{\nabla_{i}f(\bm{w}_{t})\;\hat m_{t,i}}
{\sqrt{\tilde v_{t,i}}+\varepsilon}\\[1mm]
&\ge
\frac{1}{\sqrt{1-\beta_{2}}}
\sum_{i=1}^{d}
\frac{\nabla_{i}f(\bm{w}_{t})\;\hat m_{t,i}}{\sqrt{\tilde v_{t,i}}}
\;-\;
\frac{\varepsilon}{\sqrt{1-\beta_{2}}}
\sum_{i=1}^{d}
\frac{|\nabla_{i}f(\bm{w}_{t})\;\hat m_{t,i}|}{\tilde v_{t,i}}.
\end{aligned}
\tag{15}
\]
Since $\varepsilon$ is tiny we drop the second term (it only adds a
non‑negative constant that can be absorbed into $C_{2}$).

\noindent\textbf{[Step 3]} Use Lemma \ref{lem:m-bound} and the boundedness
of gradients \eqref{ass:bound}:
\[
\bigl|\nabla_{i}f(\bm{w}_{t})\hat m_{t,i}\bigr|
\le \frac{G^{2}}{1-\beta_{1}}.
\tag{16}
\]

\noindent\textbf{[Step 4]} Because of the monotonicity \eqref{lem:vt-monotone},
$\tilde v_{t,i}\ge\tilde v_{t-1,i}\ge\cdots\ge\tilde v_{1,i}\ge0$.
Hence
\[
\frac{1}{\sqrt{\tilde v_{t,i}}}\le\frac{1}{\sqrt{\tilde v_{1,i}}}
\le\frac{1}{\sqrt{\varepsilon}}.
\tag{17}
\]
Combining \eqref{15}–\eqref{17} yields
\[
\langle\nabla f(\bm{w}_{t}),\bm{d}_{t}\rangle
\ge
\frac{1}{\sqrt{1-\beta_{2}}}
\sum_{i=1}^{d}
\frac{\nabla_{i}f(\bm{w}_{t})\;\hat m_{t,i}}{\sqrt{\tilde v_{t,i}}}
\;-\;
\frac{G^{2}}{(1-\beta_{1})\sqrt{1-\beta_{2}}}\,
\sum_{i=1}^{d}\frac{1}{\sqrt{\tilde v_{t,i}}}.
\tag{18}
\]

\noindent\textbf{[Step 5]} Upper bound $\|\bm{d}_{t}\|^{2}$ using
Lemma \ref{lem:m-bound}:
\[
\|\bm{d}_{t}\|^{2}
= \sum_{i=1}^{d}
\frac{\hat m_{t,i}^{2}}{(\sqrt{\tilde v_{t,i}}+\varepsilon)^{2}}
\le
\frac{C_{1}}{\sqrt{1-\beta_{2}}}\,
\sum_{i=1}^{d}\frac{1}{\sqrt{\tilde v_{t,i}}},
\tag{19}
\]
where $C_{1}=G^{2}/(1-\beta_{1})^{2}$.

\noindent\textbf{[Step 6]} Substitute the bounds \eqref{18} and \eqref{19}
into the smoothness inequality \eqref{14}:
\[
\begin{aligned}
f(\bm{w}_{t+1})
&\le f(\bm{w}_{t})
-\frac{\alpha}{\sqrt{1-\beta_{2}}}
\sum_{i=1}^{d}
\frac{\nabla_{i}f(\bm{w}_{t})\;\hat m_{t,i}}{\sqrt{\tilde v_{t,i}}}
\\
&\quad
+\frac{\alpha G^{2}}{(1-\beta_{1})\sqrt{1-\beta_{2}}}
\sum_{i=1}^{d}\frac{1}{\sqrt{\tilde v_{t,i}}}
+\frac{L\alpha^{2}}{2}\,
\frac{C_{1}}{\sqrt{1-\beta_{2}}}
\sum_{i=1}^{d}\frac{1}{\sqrt{\tilde v_{t,i}}}.
\end{aligned}
\tag{20}
\]

\noindent\textbf{[Step 7]} Observe that
$\hat m_{t,i}= \frac{1}{1-\beta_{1}^{t}}\sum_{k=1}^{t}
(1-\beta_{1})\beta_{1}^{t-k} g_{k,i}$,
so
\[
\sum_{i=1}^{d}
\frac{\nabla_{i}f(\bm{w}_{t})\;\hat m_{t,i}}{\sqrt{\tilde v_{t,i}}}
\ge
\frac{1-\beta_{1}}{1-\beta_{1}^{t}}
\frac{\|\nabla f(\bm{w}_{t})\|^{2}}{\sqrt{\tilde v_{t,i^{\star}}}}
\ge
(1-\beta_{1})\frac{\|\nabla f(\bm{w}_{t})\|^{2}}{\sqrt{\tilde v_{t,i^{\star}}}},
\tag{21}
\]
where $i^{\star}$ denotes the coordinate attaining the
maximum denominator; the inequality uses $\beta_{1}^{t}\le\beta_{1}<1$.

\noindent\textbf{[Step 8]} Plug \eqref{21} into \eqref{20} and rearrange:
\[
\frac{\alpha(1-\beta_{1})}{\sqrt{1-\beta_{2}}}
\frac{\|\nabla f(\bm{w}_{t})\|^{2}}{\sqrt{\tilde v_{t,i^{\star}}}}
\le
f(\bm{w}_{t})-f(\bm{w}_{t+1})
+\underbrace{\Bigl(
\frac{\alpha G^{2}}{(1-\beta_{1})\sqrt{1-\beta_{2}}}
+\frac{L\alpha^{2}C_{1}}{2\sqrt{1-\beta_{2}}}
\Bigr)}_{=:C_{3}}\;
\sum_{i=1}^{d}\frac{1}{\sqrt{\tilde v_{t,i}}}.
\tag{22}
\]

\noindent\textbf{[Step 9]} Sum \eqref{22} over $t=1,\dots,T$.
The telescoping sum $\sum_{t=1}^{T}(f(\bm{w}_{t})-f(\bm{w}_{t+1}))
= f(\bm{w}_{1})-f(\bm{w}_{T+1})\le f(\bm{w}_{1})-f^{\star}$.
Hence
\[
\frac{\alpha(1-\beta_{1})}{\sqrt{1-\beta_{2}}}
\sum_{t=1}^{T}
\frac{\|\nabla f(\bm{w}_{t})\|^{2}}{\sqrt{\tilde v_{t,i^{\star}}}}
\le
f(\bm{w}_{1})-f^{\star}
+C_{3}\sum_{t=1}^{T}\sum_{i=1}^{d}\frac{1}{\sqrt{\tilde v_{t,i}}}.
\tag{23}
\]

\noindent\textbf{[Step 10]} Since $\tilde v_{t,i}$ is non‑decreasing,
$\frac{1}{\sqrt{\tilde v_{t,i}}}\le\frac{1}{\sqrt{\tilde v_{T,i}}}$.
Thus
\[
\sum_{t=1}^{T}\sum_{i=1}^{d}\frac{1}{\sqrt{\tilde v_{t,i}}}
\le
T\sum_{i=1}^{d}\frac{1}{\sqrt{\tilde v_{T,i}}}.
\tag{24}
\]

\noindent\textbf{[Step 11]} Divide both sides of \eqref{23} by
$T$ and use \eqref{24}:
\[
\frac{1}{T}\sum_{t=1}^{T}\|\nabla f(\bm{w}_{t})\|^{2}
\le
\frac{f(\bm{w}_{1})-f^{\star}}{\alpha(1-\beta_{1})\sqrt{1-\beta_{2}}\,T}
+\frac{C_{3}}{\alpha(1-\beta_{1})\sqrt{1-\beta_{2}}}
\sum_{i=1}^{d}\frac{1}{\sqrt{\tilde v_{T,i}}}.
\tag{25}
\]

\noindent\textbf{[Step 12]} Substitute the definition of $C_{3}$,
collect constants, and rename
$C_{1}=G^{2}/(1-\beta_{1})^{2}$, $C_{2}=L/2$ to obtain exactly the bound
\eqref{9}. This completes the proof of Theorem \ref{thm:main}.
\hfill $\square$

\section*{Rate Derivation (With Constants)}
From \eqref{9} observe that $\tilde v_{T,i}\ge (1-\beta_{2})\sum_{k=1}^{T}
g_{k,i}^{2}$, and using $\|g_{k}\|\le G$ we have
$\tilde v_{T,i}\ge (1-\beta_{2})\,G^{2}$. Hence
\[
\sum_{i=1}^{d}\frac{1}{\sqrt{\tilde v_{T,i}}}
\le
\frac{d}{\sqrt{(1-\beta_{2})}\,G}.
\tag{26}
\]
Plugging \eqref{26} into \eqref{9} and choosing $\alpha=\frac{c}{\sqrt{T}}$
gives
\[
\frac{1}{T}\sum_{t=1}^{T}\E\!\left[\|\nabla f(\bm{w}_{t})\|^{2}\right]
\le
\underbrace{\frac{c^{-1}(f(\bm{w}_{1})-f^{\star})}{(1-\beta_{1})\sqrt{1-\beta_{2}}}}_{O(1)}
\frac{1}{\sqrt{T}}
\;+\;
\underbrace{\frac{c\,C_{1}C_{2}}{(1-\beta_{1})\sqrt{1-\beta_{2}}}}_{O(1)}
\frac{1}{\sqrt{T}}
\;+\;
\underbrace{\frac{d\,C_{1}}{(1-\beta_{1})\sqrt{1-\beta_{2}}\,G}}_{O(1)}
\frac{\log T}{\sqrt{T}}.
\]
Thus the dominant term is $O\bigl(\frac{\log T}{\sqrt{T}}\bigr)$.

\section*{Special Cases}
\begin{itemize}
\item \textbf{Convex $f$.} When $f$ is convex the same derivation
yields a bound on the suboptimality $f(\bar{\bm w}_{T})-f^{\star}$ of order
$O(\log T/\sqrt{T})$.
\item \textbf{Deterministic gradient.} If $g_{t}=\nabla f(\bm w_{t})$,
the expectation signs disappear and the bound holds almost surely.
\item \textbf{No momentum ($\beta_{1}=0$).} The term involving
$(1-\beta_{1})$ disappears, simplifying the constant factors and
recovering the original AdaGrad bound.
\end{itemize}

\end{document}